\section{OBJETIVOS}
\label{sec:objetivos}

\subsection{Objetivo General}
Desarrollar y validar una arquitectura de control secundario distribuido basada en Sistemas Multiagente y algoritmos de consenso en tiempo fijo (FxTS) sobre topologías de grafo, implementada mediante una plataforma de co-simulación ciber-física multitasa en hardware de borde, para garantizar la regulación de voltaje y frecuencia y el reparto proporcional de potencia activa y reactiva, en microrredes eléctricas aisladas ante contingencias en la infraestructura de comunicación.

\subsection{Objetivos Específicos}
\begin{enumerate}
    \item \textbf{Formulación Matemática y de Control (Concepción):} Formular el protocolo de control secundario distribuido en tiempo fijo fundamentado en teoría de grafos y propiedades de las matrices Laplacianas, garantizando analíticamente la convergencia de estados para la restitución del perfil de tensión nominal y la frecuencia síncrona de referencia entre los recursos energéticos distribuidos.
    
    \item \textbf{Modelado y Arquitectura de Co-Simulación (Diseño):} Diseñar un marco de co-simulación ciber-física multitasa que acople y sincronice la resolución de las dinámicas de los convertidores de potencia ($dt_{\text{micro}} = 1\text{ ms}$) con el solucionador de flujo de potencia de la red eléctrica ($dt_{\text{macro}} = 500\text{ ms}$), integrando protocolos de intercambio síncrono (ZeroMQ) y telemétrico asíncrono (MQTT).
    
    \item \textbf{Despliegue en Hardware Embebido (Implementación):} Implementar la arquitectura distribuida en un entorno experimental compuesto por dispositivos de borde (clúster Raspberry Pi 5) y contenedores modulares aislados (Docker), instanciando los recursos de generación de la configuración híbrida de validación (Diésel, Solar Fotovoltaico y BESS) sobre una red radial de ocho barras en baja tensión, e incorporando una capa de emulación de condiciones adversas de red.
    
    \item \textbf{Evaluación y Validación de Resiliencia (Evaluación):} Evaluar cuantitativamente el desempeño dinámico, el tiempo de asentamiento y la tolerancia a fallos del sistema multiagente ante escenarios de contingencia ciber-física, en particular variaciones bruscas de irradiancia y carga, pérdida de paquetes de datos y desconexión de enlaces en la topología de comunicación.
\end{enumerate}
